\section{Discussion and Limitations}\label{sec:limitations}

Within its maintained transmission model and prespecified information geometry,
the paper establishes a bounded positive implication.
The certificate maps observed Wasserstein separation into a one-sided
restriction on portfolio variance, and the zero-slack allocation rule
minimizes the resulting certified upper bound without estimating cross-asset
return covariance.
In the \ppnum[0]{n-tickers}-firm panel, only
\ppnum[2]{news-only-percentile-uniform-cap-15-pct}\%--
\ppnum[2]{news-only-percentile-concentrated-cap-15-pct}\% of portfolios across
four prespecified reference populations have conventional variance no greater
than that allocation, while equal risk weights rank materially higher under
every comparison.
The evidence is therefore an unusually low location among feasible allocations,
not covariance-free replication of the covariance-informed optimum.
Five boundaries qualify what this finding establishes and what it does not.

The first is identification.
The certificate is conditional on an antilipschitz carrier, firm-specific
slack, one coherent joint exposure law, and the return bridge.
Together, these maintained restrictions transmit observed information
separation into a statement about portfolio risk, but the text data do not
identify that transmission mechanism.
The cost is interpretive: the certificate is valid under the declared
restrictions rather than an unconditional implication of text geometry.
At the zero-slack boundary, however, the carrier constant drops out of the
normalized allocation through \Cref{cor:p3-carrier-invariance}, so the
selected risk weights depend only on the observed W$_2$ distances.
The maintained restrictions therefore determine what the low variance ranking
means economically---whether it reflects genuine risk reduction through
information separation---rather than whether the allocation exists or whether
its realized variance is low relative to the declared reference populations.
Isolating the transmission would require identified or calibrated counterparts
to the carrier, firm-specific slack, and return bridge.

The second boundary is representation dependence.
The article selection, embedding model, and output width determine the
measured distributions and may therefore alter the distances and portfolio
weights.
The representation ladder evaluates this sensitivity while holding the priced
firms and return evaluation fixed, but it does not make the canonical
Qwen3-Embedding-8B geometry invariant to measurement choices.
Two features of the ladder sharpen the interpretive consequence.
The width ladder is non-monotone: extreme compression to 64 coordinates
materially alters the point allocation and improves its in-sample variance
ranking, although the paired contrast with 256 coordinates does not survive
Holm adjustment, so the geometry is not simply degraded by truncation.
Among the matched EttaX vintage contrasts, only V3--V1 satisfies the
post-specified equivalence criterion; neither comparison involving V0 does.
The evidence therefore does not establish broad representation invariance even
within a fixed architecture.
The primary representation was prespecified, and the ladder evaluates
sensitivity rather than selecting a winner from the same return data.
Prespecified alternatives using different ground metrics, unbalanced clouds,
or held-out evaluation windows would show which features of the geometry
survive those measurement choices.

The third boundary concerns comparison-set design.
The percentile ranking is relative to four prespecified Dirichlet-capped
reference populations, and the choice of concentration parameter and cap
determines what ``low'' means.
The two concentrated populations reduce the effective-$N$ gap relative to the
uniform laws, and the ranking survives, but neither matches the candidate's
realized effective $N$.
They therefore provide concentration sensitivity rather than a
concentration-adjusted comparison, and they do not eliminate concentration as
a confound.
These four laws nevertheless remain maintained design choices.
Alternative reference populations---market-capitalization-weighted draws,
factor-tilted draws, or draws from the empirical distribution of a historical
allocation universe---could rank the candidate differently.
The evidence is that the allocation is unusually low under every declared
comparison, not under every conceivable one.

The fourth boundary is portfolio implementation.
The empirical exercise implements the standardized-asset specialization, in
which unit marginal scales make capital weights coincide with normalized risk
weights.
The general raw-capital objective in \eqref{eq:p3-certified-objective}
retains the marginal volatility scales through $A(x)^2$ and the risk-weight
mapping $q(x)$.
Constructing a capital-weight allocation from the certificate therefore
requires the same marginal-scale inputs as an inverse-volatility portfolio,
though it does not require cross-asset return covariance.
The standardized exercise isolates the certificate's information channel
cleanly, but the practical gap between the standardized and capital-weight
allocations remains unquantified.

The fifth boundary is external validity.
The empirical universe is a coverage-screened survivor panel of
\ppnum[0]{n-tickers} firms with complete return histories, not a historical
index-membership universe.
The resulting evidence characterizes firms with the required source coverage
rather than every investable firm or a changing historical universe.
The information geometry and the return evaluation also draw on the same
2018--2022 period.
The reported percentile is therefore descriptive: it is neither a
point-in-time test nor an ex-ante guarantee of future realized variance.
The expanding-cutoff sequence shows how the allocation evolves with
accumulated articles, but it does not constitute out-of-sample evidence
because the return evaluation period is unchanged.
A prospective performance claim requires a frozen point-in-time construction
in which the encoder, article corpus, and W$_2$ geometry are all fixed before
the evaluation returns are observed, together with genuinely out-of-sample
return evaluation.
The present evidence demonstrates the framework's implications for a specific
panel and representation; whether those implications survive prospective
evaluation is the immediate empirical priority.
The conclusion returns to the conditional portfolio-risk certificate and its
allocation rule within these boundaries.
